\section{The arithmetic foundation: negative MRDP}
\label{sec:foundation}

This section separates a standard deep theorem from the elementary logical
work needed to use it constructively. It also explains why importing a
classical theorem about the standard natural numbers into Lean is not the
same operation as proving the required formula in iPRA.

\subsection{The one established representation theorem we use}

Let $T_{\mathrm E}=I\Delta_0+\mathrm{EXP}$ be classical bounded arithmetic
with total exponentiation. A function-symbol presentation of elementary
arithmetic can be used instead. We need the following fixed-instance form
of the Gaifman--Dimitracopoulos theorem \cite{gd}:

\begin{theorem}[Established formalized MRDP]\label{thm:efa-mrdp}
For each fixed $\Sigma_1$ formula $\theta(e)$ of arithmetic, there is an
integer polynomial $P_0(e,y_1,\ldots,y_n)$ and a finite
$T_{\mathrm E}$ proof of
\[
 \forall e\left(\theta(e)\longleftrightarrow
       \exists y_1\cdots\exists y_n\;P_0(e,\boldsymbol y)=0\right).
\]
The witnesses range over natural numbers. The polynomial and the source
proof are fixed before $e$ varies.
\end{theorem}

The formalized result, including the strength of the source theory, is
explicitly confirmed in \cite[Section 1.2, p.~3 of the author preprint]{cmy}
and \cite[slide 4]{jerabek}. We invoke it for the one formula
$\theta(e)=\Halt(e)$ from \eqref{eq:halt-def}. Denote the right-hand
existential formula by $\Sol_0(e)$.

We are not assuming a truth predicate, an internally uniform evaluator
for all primitive recursive definitions, or an internally proved theorem
quantifying over all arithmetic formulas. The application uses one finite
source formula and one finite source proof. The effective classical MRDP
construction determines such a polynomial; listing its coefficients is
unnecessary for the logical theorem, although generating them would be
necessary for an executable end-to-end universal halting compiler.

\subsection{Negative translation into iPRA}

We give the proof transformation explicitly. For a formula $A$, its
G\"odel--Gentzen translation $A^{\mathrm N}$ is defined as follows:
\begin{align*}
 (s=t)^{\mathrm N}&=\neg\neg(s=t),&
 \bot^{\mathrm N}&=\bot,\\
 (A\land B)^{\mathrm N}&=A^{\mathrm N}\land B^{\mathrm N},&
 (A\to B)^{\mathrm N}&=A^{\mathrm N}\to B^{\mathrm N},\\
 (A\lor B)^{\mathrm N}&=\neg(\neg A^{\mathrm N}\land\neg B^{\mathrm N}),&
 (\forall x\,A)^{\mathrm N}&=\forall x\,A^{\mathrm N},\\
 (\exists x\,A)^{\mathrm N}&=\neg\forall x\,\neg A^{\mathrm N}.&&
\end{align*}
Other atomic decidable relations are first expanded into arithmetic
formulas, or treated by the same atomic double-negation clause.
Classical logical rules translate to intuitionistically valid rules. This
is verified by induction on a finite proof tree: the translated formulas
are stable, and the translated instance of excluded middle is
intuitionistically derivable. Thus only the nonlogical source axioms
require attention.

\begin{lemma}[Translation of the elementary arithmetic axioms]
Every axiom used in the translated fixed $T_{\mathrm E}$ proof has an iPRA
derivation.
\end{lemma}
\begin{proof}
Arithmetic defining equations and successor axioms are already present
in iPRA, after the standard definitional interpretation. Equality and
order are decidable, so their double negations reduce to the original
decidable statements.

Consider a fixed bounded formula $B(x,\boldsymbol a)$ occurring in an
induction axiom. Construct its primitive recursive characteristic
function $\chi_B$ by evaluating the fixed arithmetic terms and performing
the finitely bounded Boolean folds for its quantifiers. The internal
correctness proof gives
\[
 B(x,\boldsymbol a)\longleftrightarrow\chi_B(x,\boldsymbol a)=1.
\]
It also gives decidability of $B$ and equivalence
$B^{\mathrm N}\leftrightarrow B$. The translated induction axiom therefore
follows from quantifier-free induction on
$\chi_B(x,\boldsymbol a)=1$. Parameters are retained throughout.

Finally, the total-exponentiation axiom is witnessed by iPRA's primitive
recursive exponentiation function, defined by
$\operatorname{pow}(a,0)=1$ and
$\operatorname{pow}(a,Sb)=a\operatorname{pow}(a,b)$. Its agreement with the
source arithmetic graph is proved from the same recursion equations
using finite computation certificates. The negative translation of its
existential conclusion follows immediately from this stronger positive
witnessed statement.

Only finitely many bounded formulas occur in the source proof. The argument
compiles each of them separately; it makes no claim about a universal
truth evaluator for the whole language of primitive recursive arithmetic.
\end{proof}

For a $\Sigma_1$ formula with a decidable matrix, the negative translation
is intuitionistically equivalent to its double negation. For example,
if $R(e,w)$ is decidable, then
\[
 (\exists w\,R(e,w))^{\mathrm N}
 \longleftrightarrow\neg\forall w\,\neg R(e,w)
 \longleftrightarrow\neg\neg\exists w\,R(e,w).
\]
The last equivalence uses the intuitionistically valid equivalence
$\neg\exists w\,R\leftrightarrow\forall w\,\neg R$; it does not eliminate
a double negation on an existential.

Applying the translation to Theorem~\ref{thm:efa-mrdp} gives
\begin{equation}\label{eq:nn-mrdp}
 \ipra\ \vdash\quad
 \forall e\bigl(\neg\neg\Halt(e)
                \longleftrightarrow\neg\neg\Sol_0(e)\bigr).
\end{equation}
This is the exact intermediate conclusion. Claiming the positive
equivalence from this translation would be an unjustified extra step.

\begin{lemma}[Negative calibration]\label{lem:negative-calibration}
Intuitionistic propositional logic proves
\[
 (\neg\neg H\leftrightarrow\neg\neg S)
 \longrightarrow(\neg H\leftrightarrow\neg S).
\]
\end{lemma}
\begin{proof}
Assume $\neg H$ and $S$. The latter implies $\neg\neg S$, hence
$\neg\neg H$, contradicting $\neg H$. This proves $\neg H\to\neg S$.
Interchanging $H$ and $S$ proves the converse. Both arguments end in a
contradiction; neither extracts a witness from a double negation.
\end{proof}

Consequently the required internal theorem is
\begin{equation}\label{eq:negative-mrdp}
 \ipra\ \vdash\quad
 \forall e\bigl(\neg\Halt(e)\longleftrightarrow\neg\Sol_0(e)\bigr).
\end{equation}

\subsection{What this bridge proves, and what remains to formalize}

Equation~\eqref{eq:negative-mrdp} is enough for the eventual diagonal
argument. There is no need to prove the stronger constructive assertion
that every halting certificate can be transformed into a polynomial root.
The analytic part does provide positive root/nonzero-certificate maps;
the arithmetic halting/Diophantine bridge is used only negatively.

As a human mathematical proof, the argument above is an application of an
established proof-theoretic theorem followed by an explicit finitary proof
translation. It establishes existence of the required iPRA derivation.
As a Lean implementation, a source proof of the chosen fixed MRDP
instance must still be represented and checked, or the corresponding
negative iPRA derivation must be constructed directly. A parameter field
asserting \eqref{eq:negative-mrdp} is an interface for that work, not its
completion.

This treatment avoids an additional conservation theorem and avoids
asserting primitive recursive witness extraction from an arbitrary
classical MRDP proof. It also handles the natural-number witness domain
directly: Section~\ref{sec:reduction} squares the coordinates of the dense
curve, so no four-squares theorem is needed to convert natural witnesses
to integer witnesses.
